% Pseudocode for the core per-submission grading loop (README.md).
% Compile with: pdflatex pseudocode.tex
\documentclass[border=14pt]{standalone}
\usepackage[ruled,vlined,linesnumbered]{algorithm2e}
\usepackage{times}

\SetKwFunction{FFind}{find\_submissions}
\SetKwFunction{FGraded}{already\_graded}
\SetKwFunction{FExtract}{extract}
\SetKwFunction{FPrompt}{build\_user\_prompt}
\SetKwFunction{FOllama}{call\_ollama}
\SetKwFunction{FParse}{parse\_json}
\SetKwFunction{FSave}{save}
\SetKwFunction{FCsv}{write\_summary\_csv}

\SetKwInOut{Input}{Input}
\SetKwInOut{Output}{Output}

\begin{document}
\begin{algorithm}[H]
\caption{Grade submissions with a local LLM}
\Input{directory, filename\_pattern, rubric\_prompt, model, force}
\Output{one JSON + one TXT per submission, plus summary.csv}

$results \leftarrow \emptyset$\;
\ForEach{$submission \in$ \FFind{directory, filename\_pattern}}{
  \If{\FGraded{submission}\textup{ \textbf{and not} }force}{
    \textbf{continue} \tcp*{cache hit --- no Ollama call made}
  }
  $(text,\ wc,\ pc,\ ic) \leftarrow$ \FExtract{submission.path}\tcp*{local only: PyMuPDF / python-docx / openpyxl}
  $prompt \leftarrow$ \FPrompt{facts $= (wc, pc, ic)$, text $=$ text}\;
  $reply \leftarrow$ \FOllama{system $=$ rubric\_prompt, user $=$ prompt, model $=$ model,\newline
  \hspace*{4.2em} url $=$ \texttt{http://localhost:11434/api/chat}}\tcp*{the one network call, loopback only}
  $result \leftarrow$ \FParse{reply}\;
  \FSave{submission, result}\;
  $results \leftarrow results \cup \{result\}$\;
}
\FCsv{results}\;
\end{algorithm}
\end{document}
